\documentclass[a4paper]{article}
% Español (México) con XeLaTeX: fontspec para fuentes Unicode, polyglossia
% para la división silábica y los nombres del documento en la variante mexicana.
\usepackage{fontspec}
\usepackage{polyglossia}
\setdefaultlanguage[variant=mexican]{spanish}
% Los nombres chinos van como «pinyin (original)»: xeCJK da a los caracteres
% Han su propia fuente; sin ella XeLaTeX los omite con un aviso.
\usepackage{xeCJK}
\setCJKmainfont{FandolSong-Regular.otf}
% polyglossia no traduce los nombres de listings.
\AtBeginDocument{\providecommand{\lstlistingname}{}\renewcommand{\lstlistingname}{Listado}%
  \providecommand{\lstlistlistingname}{}\renewcommand{\lstlistlistingname}{Índice de listados}}
\usepackage{amsmath, amssymb}
\usepackage{graphicx}
\usepackage[margin=2.5cm]{geometry}
\usepackage[most]{tcolorbox}
\usepackage{etoolbox}
\usepackage{listings}
\usepackage{hyperref}
\usepackage{booktabs}
\usepackage{float}
\newtcolorbox{knowledgebox}[1]{enhanced, colback=blue!5!white, colframe=blue!75!black, colbacktitle=blue!75!black, coltitle=white, fonttitle=\bfseries, title={#1}, attach boxed title to top left={yshift=-2mm, xshift=2mm}, boxrule=1pt, sharp corners}
\newtcolorbox{importantbox}[1]{enhanced, colback=yellow!10!white, colframe=yellow!80!black, colbacktitle=yellow!80!black, coltitle=black, fonttitle=\bfseries, title={#1}, sharp corners}
\newtcolorbox{warningbox}[1]{enhanced, colback=red!5!white, colframe=red!75!black, colbacktitle=red!75!black, coltitle=white, fonttitle=\bfseries, title={#1}, sharp corners}
\lstset{language=Python, basicstyle=\ttfamily\small, keywordstyle=\color{blue}, stringstyle=\color{red!60!black}, commentstyle=\color{green!60!black}, breaklines=true, frame=single, numbers=left, numberstyle=\tiny\color{gray}, captionpos=b, extendedchars=true}
\newcommand{\notetitle}{Carnaval de IA Qingke (青稞): mesa redonda temática sobre LLM/MLLM}
\newcommand{\noteauthors}{Moderador: Chen Jiale (陈嘉乐) \quad Invitados: Cao Yu (曹玉), Xue Fuzhao (薛福昭), Liu Ziwei (刘子伟), Xie Tianbao (谢天宝)}
\newcommand{\notedate}{2025}
\newcommand{\videochannel}{Comunidad Qingke}
\newcommand{\videopublishdate}{2025}
\newcommand{\videoduration}{aproximadamente 99 minutos}
\newcommand{\videourl}{}
\newcommand{\videocoverpath}{cover.jpg}
\newcommand{\repourl}{https://github.com/hqhq1025/ai-course-notes}

% Footer with repo info
\usepackage{fancyhdr}
\pagestyle{fancy}
\fancyhf{}
\fancyfoot[C]{\thepage}
\fancyfoot[R]{\footnotesize\color{gray}\href{\repourl}{AI Course Notes}}
\renewcommand{\headrulewidth}{0pt}
\renewcommand{\footrulewidth}{0pt}

\begin{document}
\begin{titlepage}\centering
{\Large Notas del curso\par}\vspace{1.2cm}{\huge\bfseries \notetitle\par}\vspace{0.8cm}{\large \noteauthors\par}\vspace{0.3cm}{\large \notedate\par}\vspace{1.2cm}
\ifdefempty{\videocoverpath}{}{\includegraphics[width=0.82\textwidth,height=0.45\textheight,keepaspectratio]{\videocoverpath}\par}
\vfill
\begin{tcolorbox}[width=0.9\textwidth, colback=black!2!white, colframe=black!60, sharp corners]
\textbf{Autor o canal del video}: \videochannel\par\textbf{Fecha de publicación}: \videopublishdate\par\textbf{Duración del video}: \videoduration\par
\ifdefempty{\videourl}{}{\textbf{Enlace del video}: \href{\videourl}{\nolinkurl{\videourl}}\par}\par
\textbf{Página del proyecto}: \href{\repourl}{\nolinkurl{\repourl}}\par
\end{tcolorbox}
\end{titlepage}
\tableofcontents\newpage

%% ===== Presentación del invitado y panorama tecnológico de 2025 =====
\section{Introducción y presentación de los invitados}

Esta sesión es la mesa redonda temática sobre LLM/MLLM del Carnaval de IA Qingke (青稞 AI 嘉年华), moderada por Chen Jiale (陈嘉乐), organizada en torno a tres temas principales: \textbf{preentrenamiento (arquitectura del modelo, datos, unificación multimodal)}, \textbf{posentrenamiento (RL, Reasoning)}, y \textbf{Agentic \& Tool Use}.

\begin{knowledgebox}{Los cuatro invitados}
\begin{itemize}
  \item \textbf{Cao Yu (曹玉)}——especialista en posentrenamiento y RL de modelos multimodales grandes en cierta empresa. Con una trayectoria larga en aprendizaje por refuerzo y posentrenamiento, 2,025 fue el año en que su área «dio frutos».
  \item \textbf{Xue Fuzhao (薛福昭)}——Research Scientist en Google DeepMind, dedicado principalmente a Pretraining Scaling Law, Data-Efficient Scaling, Distillation y Model Architecture. Desde la perspectiva interna de Google, considera que el Pretraining sigue avanzando de manera constante y no se ha desacelerado.
  \item \textbf{Liu Ziwei (刘子伟)}——Universidad Tecnológica de Nanyang, Singapur. Se enfoca en tres áreas: Reasoning, Model Architecture (incluida Diffusion LM) y multimodalidad nativa, y considera que tanto en la academia como en la industria están surgiendo muchos trabajos interesantes.
  \item \textbf{Xie Tianbao (谢天宝)}——estudiante de doctorado en la Universidad de Hong Kong, enfocado principalmente en la ingeniería de Computer Use y Multimodal Agent. Con una orientación hacia Engineering Science, enfatiza el papel central de la capacidad de ingeniería en la iteración actual de los modelos grandes.
\end{itemize}
\end{knowledgebox}

\begin{importantbox}{Impresión general de 2,025: un año de IA, tres años del mundo humano}
A partir del avance de DeepSeek R1, toda la industria pisó el acelerador. Los modelos de código abierto alcanzaron rápidamente a modelos cerrados como O1/O3 en el área de Reasoning. Aplicaciones como texto a imagen y texto a video se desarrollaron con gran vitalidad. Los límites de capacidad de los modelos de vanguardia siguen expandiéndose con rapidez, y el ritmo de desarrollo está lejos de desacelerarse.
\end{importantbox}

\subsection{Resumen de la sección}
Los cuatro invitados provienen respectivamente de Google DeepMind (preentrenamiento), la academia (NTU/HKU) y la industria (posentrenamiento/RL), con perspectivas complementarias. 2,025 fue un año de desarrollo acelerado de la IA, y tanto en preentrenamiento como en posentrenamiento y en la capa de aplicación se lograron avances notables.

%% ===== Preentrenamiento: arquitectura del modelo =====
\section{Arquitectura del modelo: MoE vs Dense}

\subsection{Ventajas a corto plazo y riesgos a largo plazo de MoE}

El moderador planteó: MoE y la atención dispersa ya son prácticas dominantes; ¿la atención del grupo a la eficiencia ya superó a la atención a los resultados?

\begin{importantbox}{Xue Fuzhao (薛福昭): MoE podría ser eliminado a largo plazo}
En esencia, MoE intercambia memory (la cantidad de parámetros) por FLOPs, lo cual a su vez implica una data efficiency menor. Pero \textbf{a largo plazo los FLOPs siempre son lo más barato}: aunque la ley de Moore ya no es perfecta, la velocidad de mejora de los chips sigue siendo mucho más rápida que el crecimiento de los datos. Por lo tanto:
\begin{itemize}
  \item Todo esquema que intercambia otra cosa por FLOPs es una \textbf{solution a corto plazo}
  \item Todo esquema que intercambia FLOPs por otra cosa es una \textbf{solution a largo plazo}
\end{itemize}
El modelo podría regresar gradualmente a Dense, e incluso convertirse en Super-Dense, como el Diffusion Model.
\end{importantbox}

\textbf{Cao Yu (曹玉) (perspectiva de posentrenamiento)}: Está completamente de acuerdo en que los FLOPs son la única variable determinista. Pero añade un punto: el preentrenamiento tradicional sin datos sintéticos «posiblemente ya esté dead». Modelos como Gemini incorporan una gran cantidad de datos sintéticos desde la etapa de preentrenamiento, lo cual equivale a \textbf{intercambiar FLOPs por Data}, y luego usar Data para hacer Scaling. Si el preentrenamiento mismo ya está haciendo «FLOPs for Data», entonces MoE, como \textbf{solución de ingeniería dentro de un sistema de hardware limitado por el ancho de banda}, todavía tiene su valor.

\textbf{Liu Ziwei (刘子伟)}: Existe un Trade-off entre el corto y el largo plazo; personalmente se inclina un poco más hacia Dense Model. Pero, considerando escenarios concretos de despliegue, MoE todavía tiene sus características de uso propias. Además, hay quienes exploran la próxima generación de arquitecturas de cómputo (con un acoplamiento más estrecho entre almacenamiento y cómputo), lo cual podría influir en la elección de arquitectura dentro de 3 a 5 años.

\textbf{Xie Tianbao (谢天宝)}: La investigación en Model Architecture ya entró en una etapa de gran floración de propuestas. MoE es solo el esquema maduro de la industria; en el ámbito académico todavía hay mucha más exploración. Lo importante es que \textbf{el diseño de la arquitectura debe seguir al Learning Paradigm}: si en el futuro el Continual Learning se vuelve importante, la elección de arquitectura también cambiará en consecuencia.

\begin{warningbox}{El consenso de OpenAI: de Compute-Bound a Data-Bound}
Xue Fuzhao señaló que OpenAI declaró públicamente, en el lanzamiento de GPT-4.5, que en general ya se había pasado de Compute-Bound a Data-Bound: este es el consenso de toda la Community.
\end{warningbox}

\subsection{Resumen de la sección}
MoE es la solución de ingeniería óptima bajo las restricciones actuales de hardware, pero a largo plazo podría regresar a Dense. El paradigma de preentrenamiento ya cambió: la incorporación masiva de datos sintéticos ha difuminado los límites del «preentrenamiento tradicional». La elección de arquitectura depende, en última instancia, de la evolución del hardware y del paradigma de aprendizaje.

%% ===== Multimodalidad nativa =====
\section{Multimodalidad nativa vs multimodalidad de puente}

\subsection{Comparación de las dos rutas}

El moderador introduce la discusión: multimodalidad nativa (convierte directamente los pixel en token y, desde el inicio, se entrena con una mezcla de imagen y texto) vs multimodalidad de puente (primero se entrena el LLM y después se conecta un ViT + adapter).

\begin{importantbox}{Xue Fuzhao (薛福昭): Native es más clean, pero el Encoder tiene ventajas de ingeniería}
\textbf{Ventajas de la multimodalidad de puente (Encoder-based)}:
\begin{itemize}
  \item El Encoder y el Decoder se pueden entrenar por separado, lo que facilita la optimización
  \item Se puede comprimir de forma flexible la cantidad de Token, con mejor eficiencia de cómputo global
\end{itemize}
\textbf{Desventajas de la multimodalidad de puente}:
\begin{itemize}
  \item El objetivo de entrenamiento del Encoder (¿CLIP? ¿Detection? ¿Reconstruction?) no puede garantizar que sea completamente sin pérdida para las etapas posteriores
  \item Alta complejidad de Infra: la posición irregular de Image provoca un workload desbalanceado durante el Sequence Parallelism, lo que exige un reshuffle/recompute complejo y daña mucho el MFU
\end{itemize}
\textbf{Ventajas de Native}: el pixel token y el text token no tienen diferencia, y a nivel de Infra casi no es necesario modificar la infraestructura de un LM puro.
\end{importantbox}

\textbf{Liu Ziwei (刘子伟) (prefiere Native)}: las tres grandes ventajas de Native:
\begin{enumerate}
  \item \textbf{Consideración global de Data} — no hace falta entrenar primero el ViT y luego conectar el LM; la proporción de datos es más unificada
  \item \textbf{Emerging Ability} — desde el inicio hubo Interleaved Data, y el Early Fusion puede producir capacidades emergentes (se reporta que Gemini 3 volvió a entrenar la versión nativa)
  \item \textbf{Infra Clean} — se puede reutilizar directamente el Infra del LM, y el despliegue es más simple
\end{enumerate}
Pero por ahora falta un Codebase/Checkpoint de Native Multimodal de código abierto que todos puedan usar.

\textbf{Xie Tianbao (谢天宝)}: la Native Multimodal se discute desde Gemini 1 en 2,023. En esencia es un problema de madurez de ingeniería — la industria ya tiene las ideas, pero la capacidad de ingeniería todavía no las alcanza del todo, y sigue aprovechando los beneficios del esquema actual. \textbf{Al final, cuál esquema se queda depende de la mantenibilidad y la controlabilidad} (por analogía con la ingeniería de software).

\textbf{Cao Yu (曹玉) (perspectiva de datos)}: el desarrollo de los modelos multimodales \textbf{sigue la disponibilidad de los datos}. La multimodalidad nativa tiene mejor afinidad con los datos — para el posentrenamiento, cuanto más fuerte sea la capacidad base del base model, mejor, y la arquitectura native tiene una ventaja natural en este aspecto. Al mismo tiempo plantea una pregunta sin resolver: ¿qué modalidades debería incluir el \textbf{lado de salida} de la multimodalidad nativa? En el lado de entrada (text/vision/audio) no hay controversia, pero en el lado de salida cada quien opina distinto — los seres humanos mismos no tienen una capacidad de salida vision native.

\subsection{¿Por qué se separan LLM y VLM? ¿Cuándo se unificarán?}

El moderador señala que muchas empresas nacionales todavía publican el LLM y el VLM por separado, mientras que en el extranjero (GPT, Gemini) se tiende a modelos unificados.

\textbf{Xue Fuzhao}: los modelos de código cerrado son una caja negra; puede que internamente sea simplemente un enrutamiento entre dos modelos. Si se despliega con una estructura Encoder + Decoder, lo mejor es separarlos en dos chips — una query solo de texto (text-only) no necesita la memory del vision encoder. Ponerlos juntos complica el problema, pero a nivel de Serving es resoluble.

\textbf{Xie Tianbao}: la razón por la que se separan LLM/VLM es la misma por la que se separa el desarrollo de Coding/Math/Tool Use — desarrollar por ramas y después fusionar es una práctica común en ingeniería. Por ahora se está en un estado de «versión intermedia».

\textbf{Cao Yu}: la separación de VLM y LLM es temporal. Para una entidad operante con verdadera inteligencia, la capacidad de comprender el mundo \textbf{definitivamente no puede ser de una sola modalidad}. La separación actual se debe a que el Transformer modela discrete signal de forma mucho más madura que continuous signal. A largo plazo se superarán los problemas de ingeniería/organización/tecnología y se avanzará hacia una multimodalidad de entrada nativa.

\textbf{Liu Ziwei}: la separación a corto plazo tiene factores históricos (el Benchmark se inclina hacia la inteligencia textual, y Vision no influye mucho en la clasificación). Pero desde la perspectiva de Physical AI — que necesita decisiones de largo alcance (long-horizon action/decision) y fusión multimodal —, estos escenarios \textbf{deben unificarse}. Los seres humanos «primero tienen capacidad motora y después lenguaje», la IA es al revés, y al final también tendrá que resolver esas tareas de base que necesitan sensory input.

\subsection{Resumen de la sección}
Native vs Encoder-based cada uno tiene ventajas y desventajas; a corto plazo coexisten y a largo plazo la tendencia es hacia Native. La separación de LLM/VLM es un estado temporal de ingeniería/organización; en la era de Physical AI se unificarán necesariamente.

%% ===== Datos sintéticos y Mid-training =====
\section{Datos sintéticos y Mid-training}

\subsection{La posición de Mid-training}

\begin{importantbox}{Cao Yu (曹玉): Mid-training es el estado intermedio de entrenamiento para modelar el World Model}
La definición de Mid-training siempre ha sido poco rigurosa. Al principio se le consideraba «datos de instrucción que no alcanzan la calidad de las Golden Label de SFT». Ahora el entendimiento es más profundo:
\begin{itemize}
  \item RL Post-training es la \textbf{combinación y generalización} de conocimientos y habilidades ya existentes
  \item Estas habilidades necesitan que el modelo las \textbf{entienda y experimente} durante la etapa de Mid-training
  \item Los datos de Mid-training no solo le enseñan al modelo «cómo hacer las cosas» (Policy Model), sino que también \textbf{modelan las transiciones de estado del mundo} (World Model)
  \item Al observar datos de Mid-training como los de matemáticas: el Loss Mask es distinto al de SFT, y el modelo puede ver las transiciones de estado intermedias
\end{itemize}
En la era en que las funciones de Policy Model y World Model se comprimen en un mismo modelo, la calidad de Mid-training determina en gran medida la calidad de Post-training.
\end{importantbox}

\subsection{Los dos grandes desafíos de los datos sintéticos}

\begin{warningbox}{Xue Fuzhao (薛福昭): El riesgo central de los datos sintéticos}
\textbf{Desafío uno: la transferencia entre generaciones}: ¿puede el gain de los datos sintéticos generados por el modelo actual hacer transfer a la siguiente generación de modelos, más grandes? Intuitivamente ya es un challenge, y en la práctica también lo es. Equivale a usar el ceiling del modelo actual para entrenar a la siguiente generación.

\textbf{Desafío dos: calidad previa vs. Post-training Gain}: «reinyectar» en etapas de entrenamiento más tempranas los datos generados por el modelo de Post-training puede erosionar el gain propio de Post-training. Esto se debe a que Pre $\to$ Post es, en esencia, un proceso de mejora continua de la data quality; al introducir datos de alta calidad de manera anticipada, el rendimiento marginal de Post-training puede disminuir.

\textbf{Riesgo adicional}: los modelos grandes tienen una memoria fuerte, y los datos sintéticos inevitablemente contienen distribuciones similares a las de los datos de entrenamiento, lo que provoca un repeat implícito $\to$ overfitting $\to$ Model Collapse.
\end{warningbox}

\subsection{¿Mid-training introduce conocimiento nuevo?}

\textbf{Xie Tianbao (谢天宝)}: el consenso de la industria es que Post-training eleva el pass rate de las trajectory que ya existen en Pre/Mid-training, de pass@64 a pass@1. Si Pre/Mid-training no ha visto datos suficientemente buenos, es muy difícil que Post-training los aprenda por la fuerza. Además, la ventana de tokens de Post-training es limitada y no puede contener demasiados datos, así que hay que compensarlo en Mid-training.

\textbf{Liu Ziwei (刘子伟)}: en el ámbito académico hay trabajos del nivel de un NeurIPS Best Paper Runner-up que estudian este problema: Pre-training establece la gran distribución del conocimiento, y Mid-training empuja el boundary hacia afuera. Pero si se considera un modelo nativamente multimodal, Mid-training quizás necesite un signal distinto (no solo language), además de la capacidad de learning to learn.

\begin{knowledgebox}{Hacia un análisis científico de los datos}
Xie Tianbao señala que muchas empresas y universidades tienen una capacidad limitada de análisis de datos: la decisión de «si usar un dato y en qué etapa colocarlo» depende sobre todo del juicio humano. Conceptos como Active Learning y Self-Improvement existen, pero todavía hay una brecha grande entre ellos y la gestión manual que hacen los ingenieros al entrenar los modelos principales. Hacer más científica la gestión de datos es una dirección que vale la pena investigar a fondo.
\end{knowledgebox}

\subsection{El cuello de botella de los datos multimodales y la ruta Vision-Centric}

\textbf{Cao Yu}: los modelos de lenguaje se benefician del enorme volumen de texto de internet, pero \textbf{no puede existir un segundo internet} que ofrezca text-vision pair a la misma escala. Los animales inferiores no tienen el concepto de text y aun así sobrevivieron mucho tiempo: los VLM quizás no necesiten tantos datos para establecer una base, y deberían aprender más a través de la \textbf{interacción con el entorno} (RL o learning from experience sin reward).

\textbf{Liu Ziwei}: se puede revisar el trabajo de \textbf{autosupervisión visual} (He Kaiming (何恺明), Alusha, entre otros) de la era previa al deep learning para obtener inspiración. El equipo probó Visual Jigsaw (desordenar imágenes o videos para que el modelo recupere el orden), lo cual puede activar propiedades low-level y high-level, incluida la \textbf{inteligencia espacial}: una dirección en la que ni siquiera los modelos más fuertes actuales (GPT-5, Gemini) tienen buen desempeño.

\begin{importantbox}{Xue Fuzhao: Vision-Centric es una dirección de largo plazo}
Se opone a una ruta puramente biomimética (la vision perception de los humanos y los animales proviene en gran medida del physical touching, y es muy difícil que las máquinas obtengan ese signal). Pero \textbf{liberarse de la dependencia del texto} sí hace sense: métodos como pixel reconstruction y next frame prediction pueden proporcionar una cantidad masiva de vision token. Cuando el compute se multiplique otras 10--100 veces, los text-image pairs sin duda no darán abasto, y abordarlo directamente con vision a la fuerza es viable.
\end{importantbox}

\subsection{Resumen de la sección}
Mid-training no es solo «un SFT de calidad ligeramente menor», sino el estado intermedio de entrenamiento para modelar el World Model. Los datos sintéticos enfrentan riesgos de transferencia entre generaciones y de Model Collapse. El cuello de botella de los datos multimodales quizás pueda superarse mediante la ruta de autosupervisión Vision-Centric.

%% ===== Post-entrenamiento: RL =====
\section{Posentrenamiento y aprendizaje por refuerzo}

\subsection{Elementos centrales del RL}

\textbf{Xie Tianbao (谢天宝) (perspectiva de ingeniería)} enumeró los cuatro elementos centrales del RL:
\begin{enumerate}
  \item \textbf{Environment} — enunciado del problema (conjunto de problemas) + Sandbox (entorno de ejecución). Los requisitos difieren enormemente según el escenario; algunos incluso requieren replicar todo Internet.
  \item \textbf{Foundation Model} — la calidad del Mid-training, el Cold Start SFT y la disposición del propio modelo hacia el RL son todos importantes.
  \item \textbf{Training Infra} — el context de las tareas de RL es notablemente más largo que el de SFT; el diseño de inference engine, context management y agent memory requiere iterar de ida y vuelta.
  \item \textbf{Algoritmo de RL} — el diseño del algoritmo, a su vez, exige coordinación con environment y training infra (por ejemplo, soporte de traceback, optimización de efficiency).
\end{enumerate}
Estos cuatro elementos están altamente acoplados, implican una enorme cantidad de detalles de ingeniería, son difíciles de estudiar por completo en el ámbito académico y también es difícil dominar de una sola vez el know-how.

\begin{importantbox}{Cao Yu (曹玉): el RL es control de lazo cerrado, la definición del problema es más determinante que el algoritmo}
El RL es, en esencia, \textbf{control de lazo cerrado con señal de retroalimentación}; cada eslabón de la cadena de señal es importante. La razón por la que DeepSeek R1 tuvo éxito radica en \textbf{la claridad al elegir y definir el problema} — esto determina el resultado mucho más que el algoritmo concreto.

Los retos reales que enfrenta el RL en la actualidad:
\begin{itemize}
  \item La probabilidad de que el entorno tenga bugs es muy alta
  \item Elegir el problema es muy difícil
  \item Obtener el reward es sparse
  \item Los conflictos entre distintas tareas todavía se resuelven mediante un enfoque human-centric
\end{itemize}
El RL avanza hacia \textbf{la especialización}: quienes trabajan en coding deben convertirse en buenos reward models, y quienes trabajan en verticales (finanzas, derecho, medicina) también deben conocer el dominio.
\end{importantbox}

\subsection{La falta de Scaling Law en el RL}

\textbf{Xue Fuzhao (薛福昭)}: uno de los mayores problemas del RL es \textbf{la falta de un Scaling Law a nivel de código abierto}. Los modelos de código abierto publican versiones ya entrenadas de 7B/14B/70B, no modelos que se ubiquen en la Scaling Curve y permitan predecir por completo el siguiente tier. Esto provoca:
\begin{itemize}
  \item No saber qué ocurrirá con el modelo grande después de entrenar el modelo pequeño
  \item Que la iteración dependa del Run grande final (lento y costoso)
  \item Que el propio loss del RL suba y baje sin control, que el environment sea inestable (actualizaciones de package, problemas de asincronía) y que las conclusiones resulten muy noisy
\end{itemize}
Una vez que se establezca un Scaling Law y un Leaderboard estables, la eficiencia de iteración del RL mejorará enormemente.

\subsection{El problema de la generalización del RL}

\textbf{Cao Yu}: el RL «obtiene exactamente aquello para lo que se entrena, pero eso no significa que mejoren otros aspectos». Actualmente, la magnitud de actualización de parámetros del RL sigue estando en el rango de «pequeña perturbación» (un LoRA rank extremadamente bajo puede igualar al full-parameter), lo que provoca que la capacidad en una sola tarea difícilmente se generalice de forma natural a otros dominios. Pero, en comparación con SFT (offline imitation learning), en teoría el RL debería tener una mejor generalización — la clave es que \textbf{el Scaling del RL todavía está muy lejos de alcanzar la magnitud prometida}.

\textbf{Liu Ziwei (刘子伟)}: no se busca una fuerte generalización hacia tasks completamente unseen; basta con tener cierta \textbf{capacidad de interpolación} dentro de la matriz de tasks conocidos para que ya resulte suficientemente impresionante. Por analogía, la ola de Meta Learning de 2018 tampoco encontró una solución general de generalización. El Scaling de Task y Environment podría ser la clave de la generalización.

\textbf{Xue Fuzhao}: la definición de Generalization no es clara. Si se calcula en FLOPs (cada task muestrea 1024 traces), la generalización del RL hacia un task específico \textbf{debería ser más fuerte} (frente a SFT, que usa un task por cada sequence). El cuello de botella clave es no poder obtener environments a nivel trillion.

\begin{knowledgebox}{Liu Ziwei: distinguir lo verdadero de lo falso}
En los últimos seis meses, el RL se puso muy de moda gracias al reasoning y surgieron muchos trick de variantes. Pero trabajos recientes (como Just RL) han descubierto que muchos de esos trick no son eficaces — \textbf{el enfoque más simple y más limpio, con los hiperparámetros bien ajustados, ya alcanza lo óptimo}. Tras pasar por la bubble, la comunidad empezó a distinguir cuáles avances son reales.
\end{knowledgebox}

\subsection{Resumen de la sección}
El RL pasó de ser un simple adorno a convertirse en la herramienta central para mejorar la capacidad de los modelos, pero enfrenta tres grandes retos: la complejidad del environment, la falta de scaling law y la generalización insuficiente. El éxito de DeepSeek R1 proviene de la claridad en la definición del problema y no de la innovación algorítmica. El RL avanza hacia la especialización.

%% ===== Agentic & Tool Use =====
\section{Agentic \& Tool Use}

\subsection{Los cuellos de botella centrales de los Agent actuales}

\textbf{Xie Tianbao (谢天宝)}: El mayor avance del último año fue el paso de tareas de corta duración a tareas de varios pasos y larga duración. Pero \textbf{la calidad de las decisiones sigue siendo baja}: al observar la trajectory se nota que las decision clave a menudo fallan (el report de Deep Research no refleja los hechos, el Coding Agent introduce un bug y luego hace debug de manera repetida). El segundo problema mayor es \textbf{la gestión de Context/Memory}: cuando el multimodal agent tiene suficientes token, es difícil almacenar, gestionar y aprovechar la memory sin perder información, lo cual le impone un límite de capacidad al agent. La representación, gestión y aprovechamiento de la memory siguen siendo, en esencia, la manera de investigar de 2020--2023 (como el symbolic memory management); no han cambiado de forma fundamental.

\textbf{Liu Ziwei (刘子伟)}: Tres puntos clave:
\begin{enumerate}
  \item \textbf{Long Sequence}: entre más práctico es el Tool Use, más largo es el estado; una solución end-to-end necesita superar el action space de long sequence
  \item \textbf{Self Reflection}: la probabilidad de error en cada paso se acumula con el número de pasos; una capacidad fuerte de corrección y recuperación de errores es clave para llevarlo al real-world
  \item \textbf{Transferencia de habilidades}: si las skills que se aprenden en el Computer Use pueden transfer al Physical World (la mayoría de los robots no saben usar herramientas). La razón por la que las personas aprenden rápido a usar los íconos de la computadora es que los íconos abstraen elementos de las herramientas reales; lo inverso también podría cumplirse.
\end{enumerate}

\begin{importantbox}{Cao Yu (曹玉): El producto Agentic más Impressive de 2025}
\textbf{Software}: Claude Code, con una fuerte capacidad agentic de largo alcance, logra en el digital world lo que idealmente debería hacer una agentic AI. Esto muestra que la comprensión de Anthropic sobre Agentic \& Tool Use alcanzó un nivel muy alto. Por eso toda la industria cambió, y la mayoría de los proveedores de servicios de coding hicieron de Claude Code uno de sus servicios centrales.

\textbf{Hardware}: el teléfono Doubao (豆包手机), capaz de usar de manera automática la versión Extreme de Douyin (抖音极速版) para ganar dinero y transferirlo por WeChat, es muy impressive. Pero pocos días después se limitaron muchas funciones (no se podía abrir WeChat ni comparar precios entre plataformas), lo cual dejó al descubierto \textbf{la contradicción entre el nivel técnico de la agentic AI y las políticas de apertura de los fabricantes nacionales}.

Observación central: es difícil hacer scaling hasta trillions of environments dentro del Digital World. Un paradigm posible es \textbf{llevar el agent directamente al mundo real} (como Claude Code en software y el teléfono Doubao en hardware), dejando que el mundo mismo produzca el reward signal, en lugar de entrenar sin fin en un entorno simulado.
\end{importantbox}

\textbf{Xue Fuzhao (薛福昭)}: No es experto en agentic, pero cree en la ruta que va de Tool Use $\to$ Web Agent $\to$ \textbf{Gaming} (Minecraft $\to$ Warcraft $\to$ juegos 3A del nivel de GTA) $\to$ Physical AI. El Gaming es un paso intermedio importante: los juegos 3A ofrecen un entorno simulator muy realista.

\subsection{Resumen de la sección}
El Agent ya puede manejar tareas de larga duración, pero la calidad de las decisiones es baja y la gestión de Memory es un cuello de botella difícil. Claude Code es el producto de referencia de la agentic en 2025. El reward signal del mundo real puede ser más efectivo que el de un entorno simulado. El Gaming es el paso intermedio clave hacia la Physical AI.

%% ===== Continual Learning =====
\section{Aprendizaje continuo y Memory}

\subsection{Los retos de ingeniería del Continual Learning}

\textbf{Xie Tianbao (谢天宝)}: Tomando a Cursor como ejemplo, hace un Continual Learning simple (finalización de Type), pero el ciclo no es en tiempo real (se actualiza cada varias horas), y enfrenta problemas como messy signal, gradientes efectivos con un batch size pequeño, olvido catastrófico y ataques maliciosos. El problema central es: \textbf{antes, una vez entrenado el modelo lo dejábamos freeze para hacer inference, y nunca actualizábamos los parámetros durante la inference} —esta es una diferencia de paradigma enorme, y también una línea de investigación muy grande.

\textbf{Cao Yu (曹玉)}: Desde el punto de vista lógico, ``FLOPs for Continual Learning'' es natural y coherente —cambiar FLOPs por una mejora continua de la inteligencia. Pero el reto de ingeniería es enorme: ¿con qué método hacer el Learn (¿LoRA? ¿Sparse Update?), ¿se puede lograr un Serving Per User Per Model? Espero que en 2026 haya oportunidad de conocer qué arquitecturas de modelo o qué maneras pueden realmente lograr el Test-Time Training, y no quedarse solo en la deducción teórica y en pequeños toy example.

\subsection{Dos caminos: Parameter-based vs Context-based}

\begin{knowledgebox}{Xue Fuzhao (薛福昭): sin florituras, usar Context directamente}
El Lifelong Learning tiene dos caminos:
\begin{enumerate}
  \item \textbf{Parameter-based}—LoRA / MoE con Expert añadido / Embedding Bank. Riesgo: entre más grande el modelo, más inestable; el backpropagation con un batch size pequeño es muy risky, y la experiencia de uso a largo plazo puede degradarse.
  \item \textbf{Context-based}—armar un Prompt Database (parecido a Memory Retrieval), y escribir ahí el historial. La única desventaja es que el context es muy largo—\textbf{pero un context largo, en esencia, sigue siendo FLOPs}, lo que nos regresa al principio inicial.
\end{enumerate}
Recomendación: sin florituras, hacerlo directamente como context + retrieval. La actualización del knowledge del modelo se deja al lanzamiento periódico de nuevas versiones (de cualquier forma, ahora sale un nuevo model cada varios meses).
\end{knowledgebox}

\textbf{Liu Ziwei (刘子伟)}: Un modelo grande podría parecerse a un OS—contiene distintas regiones, donde algunos parámetros deberían actualizarse lento (Low-frequency knowledge) y otros deberían actualizarse rápido (High-frequency knowledge). Desde hace ya 20--30 años existía el concepto de Slow Weight / Fast Weight (Hinton y otros). Además, distintos dominios tienen necesidades distintas respecto al Long-tail Knowledge, y ciertas key application necesariamente deben considerar este factor a nivel de modelo o de datos. Los datos tienen por naturaleza una distribución de cola larga—la parte de baja frecuencia no se mueve y la de alta frecuencia se actualiza de manera continua—esta dynamics merece un estudio a fondo.

\subsection{Resumen de la sección}
El Continual Learning es una capacidad clave para la próxima generación de IA, pero enfrenta un triple reto de estabilidad, complejidad de ingeniería y costo de Serving. El esquema Context-based es más seguro pero necesita soporte de context largo; el esquema Parameter-based es más flexible pero con mayor riesgo.

%% ===== Panorama 2026 =====
\section{Perspectivas para 2026}

\begin{itemize}
  \item \textbf{Xue Fuzhao (薛福昭)}: espera que el preentrenamiento tenga un Paradigm Shift considerable; a nivel personal, desea aprender más conocimientos interdisciplinarios como RL.
  \item \textbf{Liu Ziwei (刘子伟)}: (1) la multimodalidad avanza hacia modelos más Unified/Native, y personas de todas las áreas —NLP, Vision, Infra— pueden contribuir; (2) cree que surgirá un gran paradigma: pasar de estar Language-Centric a Vision/Multimodal-Centric.
  \item \textbf{Cao Yu (曹玉)}: seguir impulsando el posentrenamiento con RL. En 2025, muchos compañeros que trabajaban en posentrenamiento se vieron obligados a ir y venir constantemente entre los algoritmos y la Infra, lo cual fue un proceso de crecimiento acelerado. Espera que DeepSeek no vuelva a apretar el ritmo durante el Festival de Primavera (risas). Espera hacer una investigación más pegada a la realidad práctica, y responder mediante la práctica en productos preguntas como el Continual Learning.
  \item \textbf{Xie Tianbao (谢天宝)}: consolidar sólidamente el trabajo relacionado con Multimodal Agent que tiene entre manos; experimentar a pequeña escala con Self-Improvement y Continual Learning.
\end{itemize}

\section{Síntesis y ampliación}

\begin{importantbox}{Resumen de las ideas centrales de esta sesión}
\begin{enumerate}
  \item \textbf{Arquitectura}: MoE es una solución de ingeniería a corto plazo, Dense es la dirección a largo plazo; los FLOPs siempre son lo más barato.
  \item \textbf{Multimodalidad}: Native Multimodal se ve favorable a largo plazo, la separación entre LLM y VLM es un estado temporal. En la entrada no hay controversia sobre la unificación; en la salida todavía hay discrepancias.
  \item \textbf{Datos}: los datos sintéticos enfrentan riesgos de migración entre generaciones y de Model Collapse; Mid-training es el estado intermedio de entrenamiento hacia la modelización de un World Model. El autosupervisado Vision-Centric podría resolver el cuello de botella de datos multimodales.
  \item \textbf{RL}: la naturaleza de control en lazo cerrado hace que la definición del problema sea más importante que el algoritmo; la falta de una Scaling Law, la generalización insuficiente y la complejidad del environment son los tres grandes desafíos. El RL avanza hacia la especialización.
  \item \textbf{Agent}: la baja calidad de las decisiones y la gestión de Memory son el cuello de botella central; Claude Code es el referente de 2025; llevarlo directamente al mundo real para obtener la señal de reward podría ser la vía de ruptura.
  \item \textbf{Aprendizaje continuo}: el enfoque Context-based (Prompt Database + Retrieval) es más seguro y práctico que el Parameter-based; el contexto largo sigue siendo, en esencia, una cuestión de FLOPs.
\end{enumerate}
\end{importantbox}

\subsection{Lecturas adicionales}
\begin{itemize}
  \item Informe técnico de DeepSeek-R1: un nuevo paradigma de posentrenamiento con RL
  \item Entrada de blog de OpenAI ``From Compute-Bound to Data-Bound''
  \item Trabajos relacionados con Visual Jigsaw / Spatial Intelligence Benchmark
  \item Diseño de producto de Claude Code y arquitectura Agentic
  \item Bibliografía clásica sobre Slow Weight / Fast Weight (Hinton et al.)
\end{itemize}

\end{document}
